\section{Past mass and the arbitrary-cover lower bound}
\label{hd:section}

We now compare the gauge with the actual mass of small pieces of the
trace. The lower bound uses mass accumulated before a first visit.
Section~\ref{hd:cover-section} will use successive finite tests for
the complementary upper bound. Throughout this section,
\(\Gamma=\gamma([0,\infty))\), and \(h=h_\kappa\) is the full
infimum in Equation~\eqref{gauge:definition}.
The notation \(\PP^*_{\mathrm{mix}}\) continues to mean the tilted
disk law with initial angle distributed according to \(\pi_a\).

\subsection{A spatial comparison at a stopped restart}
\label{hd:spatial}

Fix the numerical constants
\begin{equation}
 B_*=100,\qquad C_*=200,\qquad
 c_*=\frac{192}{5},\qquad C_*'=\frac{8000}{27}.
 \label{hd:koebe-constants}
\end{equation}
For \(z\in\HH\) and \(100r<R_0(z)=2\operatorname{Im}z\), let \(S\)
be the first time that \(R_S(z)=100r\) while \(z\) remains alive.
On \(\{S<\infty\}\), choose a disk map
\(F:\DD\longrightarrow H_S\) with \(F(0)=z\), normalized so that
the restarted half-plane coordinate at \(z\), after positive
rescaling, is \(e^{i\theta_S(z)}\). Thus the disk has exactly the
markings of Equation~\eqref{coef:disk-model}, and
\(|F'(0)|=100r\).
For \(K=\{|w|\le1/4\}\), the growth and distortion estimates give
\begin{align}
 \overline B(z,16r)&\subset F(K)
       \subset\overline B\left(z,\frac{400}{9}r\right)
       \subset\overline B(z,C_*r), \label{hd:koebe-sets}\\
 c_*r&\le |F'(w)|\le C_*'r\qquad(w\in K).
 \label{hd:koebe-derivatives}
\end{align}
Indeed the inner and outer growth bounds are
\(100r(1/4)/(1\pm1/4)^2\), and the derivative bounds are
\(100r(1-1/4)/(1+1/4)^3\) and
\(100r(1+1/4)/(1-1/4)^3\).
In particular, \(F(K)\) contains every closed grid square of side
\(r\) containing \(z\).

Since \(F(K)\) is compactly contained in \(H_S\), its accumulated
mass at time \(S\) is zero. Proposition~\ref{mass:restart} identifies
the actual subsequent measure there, simultaneously at all later
times, with the pushforward of the disk measure weighted by
\(|F'|^d\). Consequently, if \(T\) is the time of an additional
inner-clock increment \(0<\ell<\infty\) at \(z\), then, on
\(\{S<T<\infty\}\),
\begin{equation}
 (c_*r)^d L_\ell
 \ \le\ (\mu_T-\mu_S)(F(K))
 \ \le\ (C_*'r)^d L_\ell
 \label{hd:spatial-comparison}
\end{equation}
in these restarted coordinates. The event that this test endpoint
is reached has probability one under the corresponding finite
centered law. Conditional on the stopped past
under \(\PP_z^*\), the variable on the right is the disk-model
variable with initial angle \(\theta_S(z)\).
This is an equality of actual measures before applying distortion,
so it also permits random Borel sets and random derivative weights.
The clock increment is exactly \(\ell\): the derivative at \(z\)
cancels from ratios of conformal radii.

All uses of Equation~\eqref{hd:spatial-comparison} under a centered
law are first made on a finite stopped piece, using the raw mass
functional of Proposition~\ref{mass:construction}.
For example, if \(q<T^*=\lim_{\ell\to\infty}\tau_\ell(z)\), some
finite stopped piece extends beyond \(q\). Monotonicity therefore
bounds the mass accumulated in \(F(K)\) before \(q\) by
\((C_*'r)^dL\), with \(L\) the increasing limit in the restarted
disk model. No mass after \(T^*\) is used.

\subsection{Past mass and first visits}
\label{hd:palm-section}

The relation between natural-length-biased chordal SLE and integrated
Green-weighted two-sided radial laws was proved by
Field \cite[Theorem~1]{Field2016} for $0<\kappa\le4$ in bounded
domains with analytic boundary, and by
Zhan \cite[Theorem~4.1]{Zhan2019Decomposition} for $0<\kappa<8$.
Zhan's proof connects mass strictly after an observation time to
centered stopped laws through the conditional remaining Green potential
\cite[Proposition~2.2, equation~(2.1), and equations~(4.3)--(4.5)]{Zhan2019Decomposition}.
We use this mechanism to derive the required past-mass identity
directly for our canonical mass, from its conditional measure kernel
and finite stopped changes of law. Retaining the strict time inequality
keeps each centered-law mass evaluation before its terminal time.

Let \(Q\Subset\HH\) be a deterministic closed square and let
\(H>0\). Write \(M_q=\mu_q(Q)\), and let \(v(z)\) denote the first
visit time of \(z\). We shall use the following identity:
\begin{align}
 &\EE\int_Q
   \ind_{\{\exists q\in\mathbb Q_{\ge0}:
                 \ q<v(z),\ M_q>H\}}\,\mu_\infty(\dd z)
 \notag\\
 &\hspace{12mm}=
 \int_Q G_0(z)\,
   \PP_z^*\{\exists q\in\mathbb Q_{\ge0}:
                 \ q<T^*,\ M_q>H\}\,\dd A(z).
 \label{hd:palm}
\end{align}
The strict inequality \(q<T^*\) is part of the statement.
Here is a derivation that also specifies the measurable kernel
on its right-hand side.

First take a finite set of rational times
\(q_1<\cdots<q_N\), and set
\[
 A_j=\{M_{q_j}>H,\ M_{q_k}\le H\text{ for }k<j\}.
\]
Each \(A_j\) belongs to \(\mathcal F_{q_j}\) and does not depend on
the integration point \(z\). If a successful test in this finite
set precedes \(v(z)\), its first successful test also precedes
\(v(z)\). The first-visit identity in
Corollary~\ref{mass:first-visit} thus
identifies the corresponding random mass with
\[
 \sum_{j=1}^N
    \ind_{A_j}(\mu_\infty-\mu_{q_j})(Q).
\]
Apply Equation~\eqref{mass:kernel} with the event \(A_j\), followed
by the past-test identity in Proposition~\ref{op:potential}.
The expectation becomes
\[
 \int_Q G_0(z)
       \sum_{j=1}^N\PP_z^*(A_j,\ q_j<T^*)\,\dd A(z).
\]
These events are disjoint. If a later successful test occurs before
\(T^*\), then the first successful test does too.

For completeness, for a fixed \(A\in\mathcal F_q\) the term just
used has the measurable version
\begin{equation}
 \PP_z^*(A,q<T^*)=
 \lim_{\ell\to\infty}
 \frac{1}{G_0(z)}
 \EE\!\left[
   \ind_A G_q(z)
   \ind_{\{z\in H_q,\ s_z(q)<\ell\}}
 \right].
 \label{hd:palm-kernel}
\end{equation}
At a finite level the expectation also equals
\(\EE[\ind_A\ind_{\{q<\tau_\ell(z)<\infty\}}
G_{\tau_\ell(z)}(z)]\).
This is a stopped-past test of the finite tilted law.
The raw mass functional is fixed once, and the Loewner quantities
are jointly measurable in \(z\) and the driving path.
Thus Equation~\eqref{hd:palm-kernel} is a measurable function of
\(z\), and Tonelli's Theorem applies without a common exceptional
set for uncountably many marked points.
Finally, exhaust \(\mathbb Q_{\ge0}\) by increasing finite sets,
sorting each set before making the first-success selection.
The existence events increase even though their finite partitions
change. Monotone convergence proves Equation~\eqref{hd:palm}.

\begin{lemma}[Summable mass of bad cells]
\label{hd:bad-cells}
Fix a compact set \(E\Subset\HH\). Set
\[
 D=16(C_*')^d.
\]
For an entry \((n,i)\), write \(r=r_{ni}\), \(b=b_{ni}\), and use
the square grid of side \(r\). Call a cell \(Q\) bad when
\(\mu_\infty(Q)>2Dbr^d\).
For every sufficiently large \(n\), uniformly over all its entries
and all cells meeting \(E\),
\begin{equation}
 \EE[\mu_\infty(Q);\ Q\text{ bad}]
 \le
 4\left(\int_QG_0(z)\,\dd A(z)\right)
       \PP^*_{\mathrm{mix}}(L\ge16b).
 \label{hd:bad-mass}
\end{equation}
Moreover, almost surely, \(\mu_\infty\)-almost every point of \(E\)
belongs to no bad cell for any entry of any sufficiently late
batch.
\end{lemma}

\begin{proof}
All cells in question lie in a fixed larger compact subset of
\(\HH\) for sufficiently large \(n\).
Take \(H=Dbr^d\) in Equation~\eqref{hd:palm}.
Continuity of the increasing process \(q\mapsto\mu_q(Q)\) and
the first-visit identity show that, when \(\mu_\infty(Q)>2H\),
at least half the terminal mass of \(Q\) arrives after the
accumulated mass has exceeded \(H\).
More explicitly, if its terminal value is \(M>H\), the mass
detected by the existence event on the left of
Equation~\eqref{hd:palm} is \(M-H\): rational times detect strict
crossings, and continuity lets the accumulated mass at such
crossings decrease to \(H\).
It follows that the left side of Equation~\eqref{hd:bad-mass}
is at most twice the left side of Equation~\eqref{hd:palm}.

For \(z\in Q\), restart at conformal radius \(100r\).
Equation~\eqref{hd:koebe-sets} gives \(Q\subset F(K)\).
There is no mass in \(Q\) before the restart, and the spatial
comparison gives
\[
 \mu_q(Q)\le (C_*'r)^dL\qquad(q<T^*)
\]
in the restarted model. The time before the restart in the inner
clock is
\[
 \frac{1}{2a}\log\frac{R_0(z)}{100r}.
\]
Since \(r_{ni}\le e^{-n^2}\), this tends uniformly to infinity for
all the relevant \(z\) and entries as \(n\) increases.
Lemma~\ref{op:mixing} therefore bounds the restart-angle law by
\(2\pi_a\). The event in the right side of
Equation~\eqref{hd:palm} has probability at most
\(2\PP^*_{\mathrm{mix}}(L\ge16b)\).
Together with the first factor of two this proves
Equation~\eqref{hd:bad-mass}.

The countable collection of grid lines has zero terminal mass
almost surely, by
\(\EE\mu_\infty(\dd z)=G_0(z)\dd A(z)\).
Consequently boundaries do not affect sums over closed cells.
The integrals over cells meeting \(E\) are bounded by the integral
of \(G_0\) over a fixed larger compact. Summing
Equation~\eqref{hd:bad-mass} over every non-cleanup entry,
including all repetitions, is summable in \(n\) by
Lemma~\ref{gauge:multiplicity}.
The cleanup entries are summable as well, since
\begin{equation}
 \sum_{n=1}^{\infty}\PP^*_{\mathrm{mix}}(L\ge16n)
 =\EE^*_{\mathrm{mix}}\left\lfloor\frac{L}{16}\right\rfloor
 \le\frac{m_1}{16}.
 \label{hd:cleanup-tail}
\end{equation}
Thus the expected integral, against \(\mu_\infty|_E\), of the
number of bad-cell memberships in this entire countable family
is finite. Tonelli's Theorem implies that this number is finite
at \(\mu_\infty\)-almost every point, almost surely.
Each batch is finite, so every such point has an eventual batch
cutoff as asserted.
\end{proof}

\begin{proposition}[The lower Hausdorff bound]
\label{hd:lower}
Almost surely,
\[
 \Haus^h(\gamma([s,t]))>0
 \qquad\text{for every real }0<s<t<\infty.
\]
\end{proposition}

\begin{proof}
Apply Lemma~\ref{hd:bad-cells} on a countable compact exhaustion of
\(\HH\), and work on the resulting probability-one event together
with the conclusions about the mass in
Propositions~\ref{mass:construction} and~\ref{mass:positive}.
Fix a nontrivial interval \([u,v]\) with rational endpoints,
\(0<u<v\).
The increment \(\mu_v-\mu_u\) has positive mass on some compact
\(E\Subset\HH\) in this exhaustion, and that mass is finite.
For each $N$, let $S_N$ be $E$ minus the union of all bad cells for
entries $(n,i)$ with $n\ge N$. The countable cell families make $S_N$
Borel, and $S_N$ increases with $N$. Lemma~\ref{hd:bad-cells} says
that their union carries the increment on $E$. Choose $n_0$ with
$(\mu_v-\mu_u)(S_{n_0})>0$, and set
$\nu=(\mu_v-\mu_u)|_{S_{n_0}}$. This is a finite nonatomic measure,
carried by $\gamma([u,v])$, with $0<\nu(\HH)<\infty$.

For every entry \((n,i)\) with \(n\ge n_0\), every grid cell has
\[
 \nu(Q)\le2D b_{ni}r_{ni}^d.
\]
A bad cell has zero \(\nu\)-mass by the restriction, a good cell
has the displayed bound for \(\mu_\infty\), and a cell disjoint
from \(E\) has no \(\nu\)-mass. A set \(U\) of diameter \(s>0\)
meets at most
\(16\max\{1,(s/r_{ni})^2\}\) cells of this grid: each coordinate
projection has length at most \(s\), and it meets at most
\(s/r_{ni}+2\) grid intervals. Hence its outer measure satisfies
\begin{equation}
 \nu^*(U)\le
 32D b_{ni}r_{ni}^d
       \max\{1,(s/r_{ni})^2\}
 \qquad(n\ge n_0).
 \label{hd:retained-density}
\end{equation}

To pass to the full gauge, apply the finite-exclusion assertion of
Proposition~\ref{gauge:admissible} to
$F(s)=\sup_{\diam U\le s}\nu^*(U)$, bounded by $\nu(\HH)$.
The constant supplied by that proposition is
\[
 C_\nu=\max\left\{32D,\
       \max_{\substack{n<n_0\\1\le i\le I_n}}
          \frac{\nu(\HH)}{b_{ni}r_{ni}^d}\right\},
\]
where the maximum over an empty set is zero.
It bounds $F$ by every quadratic expression, including all the
finitely many excluded entries. Taking the full infimum therefore gives
\begin{equation}
 \nu^*(U)\le C_\nu h(\diam U).
 \label{hd:full-density}
\end{equation}
This also holds for sets of diameter zero by nonatomicity.
The argument uses finiteness of the excluded set, and does not
order its radii relative to those of later batches.

For any countable cover \(\{U_j\}\) of \(\gamma([u,v])\),
outer-measure subadditivity and Equation~\eqref{hd:full-density}
give
\[
 0<\nu(\HH)\le C_\nu\sum_j h(\diam U_j).
\]
Taking the infimum over covers with any diameter cutoff proves
\(\Haus^h(\gamma([u,v]))\ge\nu(\HH)/C_\nu>0\).
There are only countably many rational intervals. Every real
\(0<s<t\) contains a nontrivial rational closed subinterval, so
monotonicity of Hausdorff measure proves the simultaneous claim.
\end{proof}

